
\section*{Parte (b)}

\subsection*{(b-i)}

Pela definição de projeção ortogonal, temos que
\[
    P_u(v) = \alpha u.
\]

\[
    v - P_u(v) \perp u,
\]
ou seja,
\[
    \langle v - \alpha u, u \rangle = 0.
\]

Expandindo o produto interno,
\[
    \langle v, u \rangle - \alpha \langle u, u \rangle = 0.
\]

Como $\langle u, u \rangle = \|u\|^2$, segue que
\[
    \alpha = \frac{\langle v, u \rangle}{\|u\|^2}.
\]

Portanto,
\[
    P_u(v) = \frac{\langle u, v \rangle}{\|u\|^2}\, u.
\]

\subsection*{(b-ii)}
Considere a definição
\[
    \cos(u,v) := \frac{\langle u, v \rangle}{\|u\|\,\|v\|}.
\]

Pelo item (b-i), a projeção de $v$ sobre $u$ é
\[
    P_u(v) = \frac{\langle u, v\rangle}{\|u\|^2}\, u.
\]

Tomando norma em ambos os lados:
\[
    \|P_u(v)\|
    = \left\|\frac{\langle u, v\rangle}{\|u\|^2} u \right\|
    = \frac{|\langle u, v\rangle|}{\|u\|^2}\,\|u\|
    = \frac{|\langle u, v\rangle|}{\|u\|}.
\]

Dividindo por $\|v\|$, obtemos:
\[
    \frac{\|P_u(v)\|}{\|v\|}
    =
    \frac{|\langle u, v\rangle|}{\|u\|\,\|v\|}.
\]

Logo,
\[
    \cos(u,v)
    =
    \frac{\langle u, v\rangle}{\|u\|\,\|v\|}
    =
    \frac{\|P_u(v)\|}{\|v\|}
    \quad (\text{com sinal dado por } \langle u,v\rangle).
\]

Geometricamente,
\[
    \cos(u,v)
    =
    \frac{\text{cateto adjacente}}{\text{hipotenusa}},
\]
onde $\|P_u(v)\|$ é o cateto adjacente e $\|v\|$ é a hipotenusa.

\subsection*{(b-iii)}


Como os vetores têm norma $1$, vale
\[
    \cos(u,v)=\langle u,v\rangle.
\]

Calculando a norma ao quadrado da diferença:
\[
    \|u-v\|^2
    =
    \langle u-v,\,u-v\rangle.
\]

Expandindo,
\[
    \|u-v\|^2
    =
    \|u\|^2+\|v\|^2-2\langle u,v\rangle.
\]

Como $\|u\|=\|v\|=1$, obtemos
\[
    \|u-v\|^2
    =
    1+1-2\langle u,v\rangle
    =
    2-2\cos(u,v).
\]

Portanto,
\[
    \|u-v\|^2 = 2-2\cos(u,v).
\]

Agora,
\[
    \cos(u_1,v_1)\geq \cos(u_2,v_2)
\]

\[
    \iff
    2-2\cos(u_1,v_1)\leq 2-2\cos(u_2,v_2)
\]

\[
    \iff
    \|u_1-v_1\|^2 \leq \|u_2-v_2\|^2
\]

\[
    \iff
    \|u_1-v_1\| \leq \|u_2-v_2\|.
\]

Logo,
\[
    \cos(u_1,v_1) \geq \cos(u_2,v_2)
    \iff
    \|u_1 - v_1\| \leq \|u_2 - v_2\|.
\]


Para vetores de norma $1$,
\[
    \|u-v\|^2 = 2-2\cos(u,v),
\]
ou seja, o cosseno é uma transformação monótona da distância euclidiana.
Assim, quanto maior $\cos(u,v)$, menor a distância entre os vetores.
Portanto, o cosseno pode ser usado como medida de similaridade para vetores normalizados.

\subsection*{(b-iv)}

\vspace{0.5em}
\noindent A métrica da norma da diferença tem seus valores proporcionais à magnitude dos vetores avaliados. Enquanto isso, a métrica do cosseno pressupõe vetores unitários e é agnóstica à magnitude.

\noindent Assim, em um caso em que apenas a direção importa, o cosseno é preferível. Em casos em que a magnitude também é relevante, a métrica da norma é preferível ao cosseno.
